%!TEX root = ../sztuthesis_main.tex
\phantomsection
\addcontentsline{toc}{section}{Abstract}

\keywordsen{Online Examination; Remote Proctoring; C/S Architecture; Electron; Streaming Media; Information Security}
\categoryen{TP391}
\begin{abstracten}
With the popularization of online education and remote examinations, ensuring fairness and security has become a 
significant challenge. Traditional proctoring methods struggle to adapt to online scenarios, and existing remote
 monitoring tools have limitations. To address these issues, this design aims to develop a secure examination 
 information transmission and monitoring system based on the C/S architecture. The client-side, built with the 
 Electron framework, provides a graphical configuration interface, responsible for real-time capture of the 
 examinee's screen and camera feed, collecting system information (such as MAC address, IP address, CPU usage, 
 running processes, etc.), and processing/distributing the video stream locally through FFmpeg and Node-Media-Server. 
 The server-side, developed with Golang and the Gin framework, is responsible for managing client connections and 
 authentication, displaying client connection status, distributing video stream addresses, and implementing real-time bidirectional 
 communication with clients via WebSocket (such as sending commands, receiving status information, sending notifications 
 with configurable priority and recording history). The system also focuses on security, employing mechanisms such 
 as encrypted storage of local configuration, secure distribution of server access credentials, Token authentication, 
 and designing anti-cheating measures like prohibiting arbitrary exit and single-instance operation. 
 Through functional, performance, compatibility, security, and weak network tests conducted in Windows 
 and macOS environments, the results show that the system is functionally complete, performs stably, exhibits 
 good compatibility, and offers high security, effectively meeting the real-time monitoring needs for computer-based 
 exams or competitions and improving proctoring efficiency and exam fairness.
\end{abstracten}